% chapters/ch08_avionics_buses.tex — Chapter 8: Avionics Data Buses
% Scope (PLAN.md): MIL-STD-1553B (BC/RT/BM, word types, command/response),
% ARINC 429 (label/SSM/SDI/word format), Ethernet TCP vs UDP for telemetry,
% redundancy & cross-channel data link (CCDL) concepts (voting, failover),
% generic serial processing links (SPIL-type), telemetry packet structure.
% Budget 40 pp, mix 28 Q&A / 8 code / 4 concept.
% All snippets verified gcc -Wall -Wextra -std=c11 (zero warnings).

\chapter{Avionics Data Buses}
\label{ch:avionics-buses}

\begin{partnote}
\textbf{Scope} --- MIL-STD-1553B (BC/RT/BM, command/status/data words,
command--response) $\rightarrow$ ARINC 429 (label, SDI, SSM, word format)
$\rightarrow$ Ethernet TCP vs UDP for telemetry $\rightarrow$ redundancy \&
cross-channel data links (voting, failover --- the CCDL family) $\rightarrow$
generic serial processing links (SPIL-type) $\rightarrow$ telemetry packet
structure \& parsing. This chapter maps directly onto the resume:
MIL-STD-1553B and RS422 telemetry parsing at BEL, CCDL and SPIL verification at
TASL.
\end{partnote}

Where CAN (Chapter 7) is the vehicle/robotics bus, these are the
\emph{flight-grade} buses --- older, simpler, deterministic, and certified into
decades of aircraft. The recurring theme is \textbf{determinism and
redundancy}: a single bus controller commanding a fixed schedule (1553),
one-way broadcast with no contention (ARINC 429), and redundant channels that
cross-check each other (CCDL). This is the candidate's home ground.

%=====================================================================
\section{Primer}
%=====================================================================

\subsection{MIL-STD-1553B --- commanded, deterministic, redundant}

1553 is a \textbf{command/response} bus: a single \textbf{Bus Controller (BC)}
initiates every transfer; \textbf{Remote Terminals (RTs)} only speak when
commanded; a \textbf{Bus Monitor (BM)} passively records. It runs at 1\,Mbit/s,
Manchester-encoded over a dual-redundant transformer-coupled twisted pair (bus A
and bus B --- if one fails, traffic continues on the other). Communication is in
20-bit \textbf{words} (3-bit sync + 16 data + 1 parity) of three types:
\textbf{command} (from the BC: RT address, transmit/receive, subaddress, word
count), \textbf{status} (from an RT: its address plus health/condition bits),
and \textbf{data}. Because the BC owns a fixed schedule, latency and bus
behaviour are fully deterministic --- the reason 1553 has flown on military
aircraft since the 1970s.

\subsection{ARINC 429 --- one-way broadcast, label-addressed}

ARINC 429 is the civil-aircraft workhorse: a \textbf{single transmitter}
broadcasts to up to 20 receivers over a shielded twisted pair, one direction
only --- if a unit needs to talk back, it needs its own 429 bus. Each 32-bit
word is self-describing: an 8-bit \textbf{label} (octal, identifying the
parameter --- airspeed, altitude), a 2-bit \textbf{SDI} (source/destination
identifier), 19 bits of \textbf{data} (BNR binary or BCD), a 2-bit \textbf{SSM}
(sign/status matrix --- valid, no-computed-data, functional-test, sign), and an
odd-\textbf{parity} bit. Speeds are low (12.5\,kbit/s) or high (100\,kbit/s).
There's no addressing handshake and no contention because there's only ever one
talker --- determinism through simplicity.

\subsection{Ethernet, TCP vs UDP, and telemetry}

Modern avionics adds Ethernet (and the deterministic AFDX/ARINC 664 variant).
The interview-relevant distinction is \textbf{TCP vs UDP}: TCP is connection-
oriented, reliable, ordered, with retransmission and flow control --- good for
configuration and file transfer, but its retransmissions add unbounded latency.
UDP is connectionless, unordered, best-effort --- and is what \emph{telemetry}
uses, because for a stream of periodic samples a fresh sample matters more than
re-delivering a stale lost one. A dropped telemetry packet is simply superseded
by the next; you'd never want TCP stalling the stream to redeliver it. You add
your own lightweight sequence number / timestamp on top of UDP to detect loss.

\subsection{Redundancy: CCDL, voting, failover, SPIL}

Flight-critical systems run \textbf{redundant channels} (dual, triple, or quad)
computing the same control law in parallel. A \textbf{Cross-Channel Data Link
(CCDL)} lets these channels exchange their inputs and outputs every frame so
they can \textbf{vote}: with three or more channels, a \textbf{median/majority
vote} masks a single faulty channel (a divergent value is outvoted), and
\textbf{failover} drops an unhealthy channel and continues on the rest. The CCDL
must be fast, deterministic, and itself fault-tolerant, because the voting
happens every control frame. A \textbf{SPIL}-type serial processing/inter-link
is a point-to-point serial data link between cards/processors carrying such
inter-channel or inter-card data. The candidate verifies exactly these:
CCDL integrity/timing/failover and SPIL packet-level behaviour.

%=====================================================================
\section{Q\&A Bank}
%=====================================================================

\tierbanner{Tier 1 --- Screening (phone / HR-technical)}%
{Two to four sentences. Avionics-bus vocabulary --- expected for any
flight-systems role.}

\question{Q1.1}{What kind of bus is MIL-STD-1553B?}
A deterministic command/response military avionics bus at 1\,Mbit/s: a single
Bus Controller commands all traffic, Remote Terminals respond only when
addressed, and the dual-redundant transformer-coupled bus tolerates a cable
fault. It's been standard on military aircraft since the 1970s.

\question{Q1.2}{What are the three terminal types on a 1553 bus?}
Bus Controller (BC) --- initiates every transfer and owns the schedule; Remote
Terminal (RT) --- a subsystem that responds only when commanded; Bus Monitor
(BM) --- a passive listener that records traffic without participating.

\question{Q1.3}{What are the three 1553 word types?}
Command word (BC$\rightarrow$RT: RT address, transmit/receive bit, subaddress,
word count), status word (RT$\rightarrow$BC: RT address plus condition/health
bits), and data word (16 bits of payload). Each is 20 bits: 3 sync + 16 + 1
parity.

\question{Q1.4}{What kind of bus is ARINC 429?}
A civil-avionics one-way broadcast bus: a single transmitter sends 32-bit
labelled words to up to 20 receivers over a twisted pair, one direction only.
No addressing handshake, no contention --- simplicity gives determinism.

\question{Q1.5}{What's in an ARINC 429 word?}
An 8-bit label (the parameter ID, in octal), 2-bit SDI (source/destination),
19 data bits (BNR or BCD), 2-bit SSM (sign/status), and an odd-parity bit ---
32 bits total, self-describing.

\question{Q1.6}{Why does telemetry use UDP rather than TCP?}
Telemetry is a stream of periodic samples where a fresh value supersedes a lost
one, so UDP's best-effort, low-latency delivery fits; TCP's retransmission and
ordering would stall the stream to re-deliver stale data. You add your own
sequence numbers to detect loss.

\question{Q1.7}{What is a Cross-Channel Data Link (CCDL)?}
A link between redundant flight-control channels that lets them exchange their
inputs and outputs each frame, so they can cross-check and vote --- the basis of
masking a faulty channel and failing over. It must be fast and deterministic.

\question{Q1.8}{What is redundancy voting?}
With multiple channels computing the same result, a vote (majority for
discretes, median for analog values) selects the agreed output and masks a
single divergent channel --- so one fault doesn't propagate to the actuator.

\question{Q1.9}{Why is 1553 dual-redundant?}
It has two independent buses (A and B); if a transceiver, stub, or cable fails
on one, the BC retries on the other, so a single physical fault doesn't lose
communication --- essential for flight-critical subsystems.

\question{Q1.10}{What is Manchester encoding and why does 1553 use it?}
Each bit is a mid-bit transition (the encoding carries its own clock), so the
receiver recovers timing from the data and there's no DC component --- ideal for
transformer coupling. It costs double the bandwidth but gives robust,
self-clocking, galvanically-isolated signalling.

\tierbanner{Tier 2 --- Core technical (panel round)}%
{A short paragraph.}

\question{Q2.1}{Walk a 1553 BC-to-RT (receive) transfer.}
The BC sends a \textbf{command word} (RT address, T/R = receive, subaddress,
word count $N$), immediately followed by $N$ \textbf{data words}. The addressed
RT validates the command, receives the data, and after the required response
time replies with its \textbf{status word} (its address + condition bits). For
an RT-to-BC (transmit) transfer, the BC sends the command (T/R = transmit) and
the RT responds with its status word \emph{then} the requested data words. Every
exchange is BC-initiated and bounded in time --- the determinism that makes 1553
analysable.

\question{Q2.2}{Decode a 1553 command word's fields.}
The 16 data bits are: RT address (5 bits, 0--31; 31 is broadcast), T/R bit
(1 = transmit from RT, 0 = receive to RT), subaddress (5 bits --- selects the
data buffer/function; subaddresses 0 and 31 signal mode codes), and word count
(5 bits, where 0 means 32 words). So you can read off ``RT 5, receive, subaddress
1, 2 words'' directly from the bit fields (Coding Gym).

\question{Q2.3}{What does the 1553 status word tell the BC?}
The RT's address (echo, for validation) plus condition bits: Message Error,
Instrumentation, Service Request, Broadcast Command Received, Busy, Subsystem
Flag, Dynamic Bus Control Acceptance, and Terminal Flag. The BC inspects these
to confirm the RT received the message cleanly and is healthy; a set Message
Error or Terminal Flag triggers retry or fault handling. It's the RT's only way
to report status, since RTs never speak unbidden.

\question{Q2.4}{Decode an ARINC 429 word, field by field.}
From the 32-bit word: bits 1--8 the \textbf{label} (octal parameter ID, and note
the label is bit-\emph{reversed} on the wire), bits 9--10 the \textbf{SDI}
(which source/sink, or part of the data), bits 11--29 the \textbf{data} (BNR
two's-complement or BCD), bits 30--31 the \textbf{SSM} (Normal Operation /
No-Computed-Data / Functional-Test / Failure-Warning, plus sign for BNR), bit 32
\textbf{parity} (odd). The receiver checks parity, reads the SSM to know if the
data is valid, then scales the data per the label's definition.

\question{Q2.5}{What is the SSM and why must you check it before using the
data?}
The Sign/Status Matrix reports the \emph{validity and sign} of the word: it can
say Normal Operation, No-Computed-Data (the source can't compute this parameter
right now), Functional-Test, or Failure-Warning. If you scale and use the data
bits without checking the SSM, you can act on a number the source has flagged as
invalid or test data --- a serious integration bug. Always gate on SSM = Normal
before trusting the value.

\question{Q2.6}{Why is ARINC 429's label transmitted bit-reversed, and what's
the practical consequence?}
Historically the label is transmitted with its bits in reverse order relative to
the rest of the word (a legacy of the standard). The practical consequence is
that firmware/test code must reverse the 8 label bits to recover the octal label
the documentation lists --- forget the reversal and every label looks wrong, a
classic ``the bus decode is garbage'' bug when first bringing up a 429 receiver.
A simple 8-bit reverse (Coding Gym) fixes it.

\question{Q2.7}{TCP vs UDP --- when would an avionics system actually use TCP?}
For \emph{reliable, ordered, non-real-time} transfers: uploading a
configuration or software update, downloading a recorded data file, or a
maintenance/diagnostic session where every byte must arrive correctly and
latency doesn't matter. The streaming, periodic, latency-sensitive traffic
(telemetry, sensor data) uses UDP (or a deterministic protocol like AFDX). The
rule: reliability-over-latency $\rightarrow$ TCP; freshness-over-reliability
$\rightarrow$ UDP.

\question{Q2.8}{How does median voting across three channels mask a fault?}
Each of three redundant channels computes the same value; the voter takes the
\textbf{median} (middle value). If one channel is faulty and produces a wild
value, it sorts to the top or bottom and the median is one of the two good,
agreeing channels --- so the fault is masked without needing to know \emph{which}
channel failed. Median (not average) is used because an average would be dragged
toward the bad value; the median rejects the outlier entirely. This runs every
control frame (Coding Gym).

\question{Q2.9}{What does a CCDL actually carry, and why every frame?}
It carries each channel's sensor inputs and/or computed outputs to the other
channels so that, within the same control frame, every channel has all channels'
data and can vote on a common result before driving the actuator. It must happen
\emph{every} frame because the vote must use current data --- stale cross-channel
data would let channels diverge. Hence the CCDL's defining requirements:
bounded latency, frame-deterministic delivery, and its own integrity checking
(the things the resume says she verifies: data integrity, timing, failover).

\question{Q2.10}{How would you structure and parse a telemetry frame robustly?}
A frame has a \textbf{sync word} (a fixed pattern to find frame start), a
\textbf{length} field, the \textbf{payload}, and a \textbf{CRC/checksum}. The
parser scans the byte stream for the sync pattern, reads the length, validates
the checksum over the frame, and only then delivers the payload; on a checksum
failure it discards and \emph{re-scans for the next sync} rather than trusting
byte alignment. This tolerates dropped/inserted bytes and noise --- the same
``never trust alignment, validate every frame'' discipline as the UART and CAN
chapters, and exactly the RS422/1553B telemetry parsing on the resume.

\question{Q2.11}{What is a SPIL-type link in this context?}
A serial processing/inter-processor link --- a point-to-point serial data link
between cards or processors carrying inter-channel or inter-card data (here, over
RS422). You verify it at the \emph{packet level}: framing, integrity, timing,
throughput, and behaviour under fault/loss. It's conceptually a private,
deterministic serial pipe between specific endpoints, distinct from a shared bus
like 1553 or CAN, and is validated with the same framing/CRC/timing rigour.

\tierbanner{Tier 3 --- Trap \& depth (principal-engineer level)}%
{A paragraph plus the trap. The defensible depth behind 1553B / CCDL / SPIL on
the resume.}

\question{Q3.1}{Why does median voting need an \emph{odd} number of channels,
and what happens with only two?}
Median/majority voting needs a tie-breaker: with three channels a single
outlier is outvoted by the two agreeing ones, so the fault is both \emph{masked}
and implicitly \emph{identified} (the disagreeing channel). With only \textbf{two}
channels you can \emph{detect} a disagreement but not resolve it --- you don't
know which one is right, so dual systems can only fail-safe (e.g.\ disengage,
revert to a default, or use a separate monitor), not fail-operational. \trap the
trap is assuming ``redundancy = keeps working''; dual redundancy gives
fail-\emph{detection}, triple gives fail-\emph{operational} (survive one fault
and keep flying), quad survives two. Knowing that voting topology determines
\emph{fail-safe vs fail-operational} is core flight-controls judgement and
directly frames CCDL verification.

\question{Q3.2}{Two redundant channels must agree to vote --- but their clocks
and sensors differ slightly. How do you vote without false fault trips?}
Real channels are never bit-identical: sensors have noise/tolerance, ADCs
quantise differently, and clocks drift, so raw values differ even when all are
healthy. The voter must use a \textbf{threshold/tolerance band} (and often
temporal filtering --- a channel must disagree beyond the band for several
consecutive frames) before declaring a fault, and channels must be
\textbf{frame-synchronised} via the CCDL so they're comparing data from the same
sample instant. \trap the trap is an exact-equality vote, which trips constantly
on normal noise (nuisance failover), or an un-synchronised vote that compares
this-frame data against last-frame data (phase error read as disagreement). The
real design tunes the disagreement threshold and the persistence count, and
keeps the channels frame-locked through the CCDL --- precisely the
``integrity/timing/failover'' verification on the resume.

\question{Q3.3}{An RT intermittently sets the Message Error bit on a 1553 bus.
Walk the investigation.}
Message Error means the RT detected a problem with the command/data it received
(bad parity, wrong word count, gap/sync error, or an illegal command). Hunt:
(1) \textbf{physical} --- scope the Manchester waveform on both buses for
amplitude/reflection issues (stub length, coupling transformer, termination),
since 1553 is unforgiving of stub-length and impedance errors; (2)
\textbf{timing} --- check the BC's inter-message gaps and the RT's response-time
window; (3) \textbf{protocol} --- is the BC sending an illegal subaddress/word
count that this RT rejects; (4) \textbf{redundancy} --- does it happen on bus A,
B, or both (both $\rightarrow$ BC/RT logic; one $\rightarrow$ that channel's
hardware). \trap the trap is treating Message Error as ``the RT is broken'';
it's the RT \emph{reporting} a received-message fault, so the cause is usually on
the BC side or the physical layer. The bus-monitor capture plus a scope on the
coupled waveform localises it --- the same instrument-led method across the
book.

\question{Q3.4}{Why can't you just cast a received 1553/ARINC/telemetry buffer
onto a C struct, and what's the correct parser?}
Three compounding hazards (Chapters 1 \& 3): \textbf{endianness} (avionics words
are often big-endian / specific bit orders, the host is usually little-endian),
\textbf{bit-field layout} (implementation-defined, so a bitfield struct over a
429 word or 1553 word may map differently per compiler), and \textbf{the
wire's own quirks} (the ARINC label is bit-reversed; 1553 packs sub-16-bit
fields). The correct parser reads the raw words and extracts each field with
explicit shifts and masks (and explicit byte/bit-order conversion), validates
parity/CRC and the status/SSM, and never trusts struct overlay. \trap the trap
is the convenient \code{(word\_t*)buf} cast that ``works'' on one build and
silently corrupts on another compiler or a big-endian target --- which is
exactly why the resume's 1553B/RS422 parsing is done at the bit/byte level, the
field-by-field discipline this whole book reinforces.

\question{Q3.5}{Compare CAN, 1553, and ARINC 429 on determinism --- when is
each the right choice?}
\textbf{1553}: fully deterministic by central schedule (the BC dictates every
transfer), dual-redundant, certified --- the choice for flight-critical
command/response on military platforms, at the cost of a heavy BC and 1\,Mbit/s
ceiling. \textbf{ARINC 429}: deterministic by simplicity (one talker, no
contention), cheap and robust, but one-way and low-rate --- ideal for
broadcasting a sensor's parameters to many listeners on civil aircraft.
\textbf{CAN}: deterministic by \emph{priority} (lowest ID wins) but multi-master
and event-capable --- flexible and cheap, great for vehicles/UAV subsystems, but
worst-case latency needs a bus-load analysis (Chapter 7) and it lacks 1553's
central schedule and certification pedigree. \trap the trap is ranking them
``best to worst''; they encode \emph{different determinism philosophies}
(scheduled vs broadcast vs prioritised) for different criticality and topology.
Articulating that --- and that a UAV may use \emph{all three} (CAN for servos,
429-style for sensors, a 1553/CCDL core for flight-critical voting) --- is the
systems-integration view the resume points to.

%=====================================================================
\section{Coding Gym}
%=====================================================================

\begin{partnote}
Eight host-compilable problems (verified under
\code{gcc -Wall -Wextra -std=c11}) implementing the field extraction, parity,
voting, and framing you actually perform validating 1553B / ARINC 429 / CCDL /
telemetry --- the byte-level parsing the resume describes.
\end{partnote}

%---------------------------------------------------------------
\begin{gymproblem}{Pack and extract an ARINC 429 word}

\textbf{Problem.} Assemble a 32-bit ARINC 429 word from label/SDI/data/SSM and
extract the fields back.

\textbf{Hints.}
\begin{itemize}
\item Label bits 1--8, SDI 9--10, data 11--29, SSM 30--31, parity 32.
\end{itemize}

\attempt

\textbf{Solution.}
\begin{cgym}
#include <stdint.h>
#include <stdio.h>

static uint32_t a429_pack(uint8_t label, uint8_t sdi, uint32_t data19, uint8_t ssm) {
    uint32_t w = 0;
    w |= label;                          /* bits 1-8   */
    w |= (uint32_t)(sdi & 0x3u) << 8;    /* bits 9-10  */
    w |= (data19 & 0x7FFFFu) << 10;      /* bits 11-29 */
    w |= (uint32_t)(ssm & 0x3u) << 29;   /* bits 30-31 */
    return w;
}
static uint8_t  a429_label(uint32_t w) { return (uint8_t)(w & 0xFFu); }
static uint8_t  a429_sdi(uint32_t w)   { return (uint8_t)((w >> 8) & 0x3u); }
static uint32_t a429_data(uint32_t w)  { return (w >> 10) & 0x7FFFFu; }
static uint8_t  a429_ssm(uint32_t w)   { return (uint8_t)((w >> 29) & 0x3u); }

int main(void) {
    uint32_t w = a429_pack(0312u & 0xFF, 0x1u, 0x5A5Au, 0x3u);
    printf("label=0%o sdi=%u data=0x%05X ssm=%u\n",
           a429_label(w), a429_sdi(w), a429_data(w), a429_ssm(w));
    /* label=0312 sdi=1 data=0x05A5A ssm=3 */
    return 0;
}
\end{cgym}

Each field is extracted with an explicit shift and mask --- \emph{not} a bitfield
struct overlay, which would be non-portable across compilers and endianness
(Q3.4). The label prints in octal because that's how ARINC documents it. This is
the literal core of a 429 receive driver and of the kind of telemetry parser the
resume describes.

\followups
\begin{enumerate}
\item \emph{``Where's the parity bit?''} $\rightarrow$ Bit 32 (next problem);
  set for odd parity over the lower 31 bits.
\item \emph{``Scale BNR data.''} $\rightarrow$ The data field is
  two's-complement with a per-label LSB weight; multiply by the resolution.
\item \emph{``Why mask each field on insert?''} $\rightarrow$ Defensive --- a
  too-wide argument would corrupt the neighbouring field.
\end{enumerate}
\end{gymproblem}

%---------------------------------------------------------------
\begin{gymproblem}{ARINC 429 odd parity}

\textbf{Problem.} Set bit 32 for odd parity and provide a check.

\textbf{Hints.}
\begin{itemize}
\item Count the 1s in bits 1--31; if even, set bit 32 to make the total odd.
\end{itemize}

\attempt

\textbf{Solution.}
\begin{cgym}
#include <stdint.h>
#include <stdio.h>

static int popcount32(uint32_t x) { int n = 0; while (x) { n++; x &= x - 1; } return n; }

static uint32_t a429_set_parity(uint32_t w) {
    if ((popcount32(w & 0x7FFFFFFFu) & 1) == 0)   /* lower 31 bits are even? */
        w |= 0x80000000u;                          /* set bit 32 -> total odd */
    return w;
}
static int a429_parity_ok(uint32_t w) { return popcount32(w) & 1; }   /* odd == good */

int main(void) {
    uint32_t w = a429_set_parity(0x0A5A5u);
    printf("parity_ok=%d\n", a429_parity_ok(w));   /* 1 */
    return 0;
}
\end{cgym}

ARINC 429 uses \textbf{odd} parity over the whole word, so a healthy received
word always has an odd number of 1s. The receiver recomputes and rejects any
word with even parity --- the first integrity gate before even looking at the
SSM. It catches single-bit errors (the dominant noise failure) but, like UART
parity, misses even-count errors --- which is why critical data also carries
application-level validity (the SSM) and often a higher-layer checksum.

\followups
\begin{enumerate}
\item \emph{``Why odd, not even?''} $\rightarrow$ Odd parity guarantees at least
  one 1, so an all-zero (dead) line is detectably invalid.
\item \emph{``Catch a 2-bit error?''} $\rightarrow$ Parity can't; rely on SSM /
  range checks / higher-layer CRC.
\item \emph{``Where does parity sit relative to SSM in validation?''}
  $\rightarrow$ Parity first (word intact?), then SSM (data valid?).
\end{enumerate}
\end{gymproblem}

%---------------------------------------------------------------
\begin{gymproblem}{Reverse the ARINC 429 label}

\textbf{Problem.} The label is transmitted bit-reversed; reverse 8 bits to
recover the documented octal label.

\textbf{Hints.}
\begin{itemize}
\item Swap nibbles, then pairs, then adjacent bits.
\end{itemize}

\attempt

\textbf{Solution.}
\begin{cgym}
#include <stdint.h>
#include <stdio.h>

static uint8_t reverse8(uint8_t b) {
    b = (uint8_t)((b >> 4) | (b << 4));                       /* swap nibbles */
    b = (uint8_t)(((b & 0xCCu) >> 2) | ((b & 0x33u) << 2));   /* swap pairs   */
    b = (uint8_t)(((b & 0xAAu) >> 1) | ((b & 0x55u) << 1));   /* swap bits    */
    return b;
}

int main(void) {
    printf("reverse8(0x0A)=0x%02X\n", reverse8(0x0A));   /* 0x50 */
    return 0;
}
\end{cgym}

This $O(1)$ bit-reverse is the fix for the most common ARINC 429 bring-up bug:
the label on the wire is reversed relative to the documented octal value, so
without this every parameter looks like the wrong label and the whole decode
appears broken (Q2.6). It's the same mask-and-swap pattern as the reverse-bits
problem in Chapter 1, applied to a real protocol quirk.

\followups
\begin{enumerate}
\item \emph{``Why does the standard reverse it?''} $\rightarrow$ Historical wire
  convention; you just have to honour it.
\item \emph{``Table-driven reverse for speed?''} $\rightarrow$ A 256-entry
  lookup if labels are decoded at high rate.
\item \emph{``How would you catch this bug on a logic analyser?''}
  $\rightarrow$ Known label (e.g.\ 0310 altitude) decodes wrong until you
  reverse.
\end{enumerate}
\end{gymproblem}

%---------------------------------------------------------------
\begin{gymproblem}{Decode a MIL-STD-1553 command word}

\textbf{Problem.} Extract RT address, T/R, subaddress, and word count from a
1553 command word.

\textbf{Hints.}
\begin{itemize}
\item RT addr bits 15--11, T/R bit 10, subaddress 9--5, word count 4--0
  (0 means 32).
\end{itemize}

\attempt

\textbf{Solution.}
\begin{cgym}
#include <stdint.h>
#include <stdio.h>

typedef struct { int rt, tr, sa, wc; } cmd1553_t;

static cmd1553_t cmd1553_decode(uint16_t w) {
    cmd1553_t c;
    c.rt = (w >> 11) & 0x1F;   /* RT address  */
    c.tr = (w >> 10) & 1;      /* 1=transmit  */
    c.sa = (w >> 5)  & 0x1F;   /* subaddress  */
    c.wc =  w        & 0x1F;   /* word count (0 => 32) */
    return c;
}

int main(void) {
    cmd1553_t c = cmd1553_decode(0x2C22);
    printf("rt=%d tr=%d sa=%d wc=%d\n", c.rt, c.tr, c.sa, c.wc);  /* 5 1 1 2 */
    return 0;
}
\end{cgym}

The command word is the BC's instruction: ``RT 5, transmit, subaddress 1, 2
words.'' Subaddresses 0 and 31 are special (they signal \emph{mode commands}
like synchronise or transmit-status, with the word-count field becoming a mode
code), and word count 0 means 32 words --- two edge cases a correct decoder must
handle. A bus monitor decodes exactly this to follow the traffic.

\followups
\begin{enumerate}
\item \emph{``What does subaddress 0 or 31 mean?''} $\rightarrow$ A mode command;
  the WC field is then a mode code, not a count.
\item \emph{``RT address 31?''} $\rightarrow$ Broadcast to all RTs (no status
  response expected).
\item \emph{``Word count 0?''} $\rightarrow$ 32 data words (the field can't
  encode 32 in 5 bits).
\end{enumerate}
\end{gymproblem}

%---------------------------------------------------------------
\begin{gymproblem}{Decode a MIL-STD-1553 status word}

\textbf{Problem.} Extract the RT address and the key condition bits from a
status word.

\textbf{Hints.}
\begin{itemize}
\item RT addr 15--11, Message Error bit 10, Service Request bit 8, Busy bit 3.
\end{itemize}

\attempt

\textbf{Solution.}
\begin{cgym}
#include <stdint.h>
#include <stdio.h>

typedef struct { int rt, msg_err, sreq, busy; } stat1553_t;

static stat1553_t stat1553_decode(uint16_t w) {
    stat1553_t s;
    s.rt      = (w >> 11) & 0x1F;
    s.msg_err = (w >> 10) & 1;    /* Message Error  */
    s.sreq    = (w >> 8)  & 1;    /* Service Request */
    s.busy    = (w >> 3)  & 1;    /* Busy            */
    return s;
}

int main(void) {
    stat1553_t s = stat1553_decode(0x2800);   /* RT 5, all clear */
    printf("rt=%d msg_err=%d sreq=%d busy=%d\n", s.rt, s.msg_err, s.sreq, s.busy);
    return 0;
}
\end{cgym}

The status word is the RT's only voice: the BC reads it after every transfer to
confirm health. Message Error means the RT rejected the command/data (Q3.3);
Busy means it can't service the request now; Service Request asks the BC for
attention. The BC's fault logic keys off these bits --- a set Message Error or
Terminal Flag drives a retry or a declared RT fault.

\followups
\begin{enumerate}
\item \emph{``Why does the RT echo its own address in the status word?''}
  $\rightarrow$ So the BC can confirm the \emph{correct} RT responded.
\item \emph{``RT doesn't respond at all --- BC action?''} $\rightarrow$
  Response-timeout $\rightarrow$ retry on the other bus, then declare the RT
  failed.
\item \emph{``Broadcast command --- is there a status word?''} $\rightarrow$ No
  status response to a broadcast (RT 31).
\end{enumerate}
\end{gymproblem}

%---------------------------------------------------------------
\begin{gymproblem}{Find and validate a telemetry frame}

\textbf{Problem.} Scan a byte stream for a sync pattern and validate the frame's
checksum.

\textbf{Hints.}
\begin{itemize}
\item Find sync (e.g.\ \code{0xEB 0x90}); sum the frame bytes; compare to the
  trailing checksum.
\end{itemize}

\attempt

\textbf{Solution.}
\begin{cgym}
#include <stdint.h>
#include <stdio.h>

static int tlm_find(const uint8_t *buf, int n, int *start) {
    for (int i = 0; i + 1 < n; i++)
        if (buf[i] == 0xEB && buf[i + 1] == 0x90) { *start = i; return 1; }
    return 0;
}
static int tlm_validate(const uint8_t *f, int len) {
    uint8_t sum = 0;
    for (int i = 0; i < len - 1; i++) sum = (uint8_t)(sum + f[i]);
    return sum == f[len - 1];          /* trailing byte is the checksum */
}

int main(void) {
    uint8_t buf[] = {0x00, 0xEB,0x90, 0x03, 0x10,0x20, 0xAE};
    int st;
    tlm_find(buf, 7, &st);
    printf("found@%d valid=%d\n", st, tlm_validate(buf + st, 6));   /* 1 1 */
    return 0;
}
\end{cgym}

The parser \emph{re-finds the sync} in the stream rather than assuming byte
alignment, then validates the checksum before trusting the payload --- so a
dropped or inserted byte (noise on the RS422 line) is detected and the parser
resynchronises on the next sync word instead of mis-decoding every following
field. This is the robust framing the resume's RS422/telemetry parsing relies
on (Q2.10).

\followups
\begin{enumerate}
\item \emph{``Replace the sum with a CRC.''} $\rightarrow$ Stronger burst-error
  detection (Chapter 13); same find-then-validate structure.
\item \emph{``Two sync bytes appear inside the payload --- problem?''}
  $\rightarrow$ Possible false sync; the length+checksum reject it, then
  re-scan.
\item \emph{``Stream split across reads (DMA buffer wrap)?''} $\rightarrow$
  Carry a small residual between buffers; never assume a frame fits one read.
\end{enumerate}
\end{gymproblem}

%---------------------------------------------------------------
\begin{gymproblem}{CCDL median-of-three voting}

\textbf{Problem.} Vote three redundant channel values by selecting the median.

\textbf{Hints.}
\begin{itemize}
\item The median is the value that lies between the other two.
\end{itemize}

\attempt

\textbf{Solution.}
\begin{cgym}
#include <stdint.h>
#include <stdio.h>

static int32_t median3(int32_t a, int32_t b, int32_t c) {
    if ((a >= b && a <= c) || (a <= b && a >= c)) return a;
    if ((b >= a && b <= c) || (b <= a && b >= c)) return b;
    return c;
}

int main(void) {
    printf("%d %d\n", median3(10, 30, 20), median3(5, 5, 9));   /* 20  5 */
    return 0;
}
\end{cgym}

Median voting masks a single faulty channel: if one channel returns a wild value
it sorts to an extreme and the median is one of the two agreeing channels
(Q2.8). Median --- not average --- is essential, because an average is dragged
toward the bad value while the median rejects the outlier entirely. In a real
voter you'd pair this with a tolerance band and a persistence count so sensor
noise doesn't trip a false failover (Q3.2). This is the arithmetic heart of
CCDL voting.

\followups
\begin{enumerate}
\item \emph{``Add a disagreement threshold.''} $\rightarrow$ Flag a fault only
  if the outlier exceeds the median by more than a band for $N$ frames.
\item \emph{``Two channels only?''} $\rightarrow$ Can detect but not vote;
  fail-safe instead (Q3.1).
\item \emph{``Four channels?''} $\rightarrow$ Average the middle two, or
  drop-the-outlier then vote --- survives two faults.
\end{enumerate}
\end{gymproblem}

%---------------------------------------------------------------
\begin{gymproblem}{Failover channel selection}

\textbf{Problem.} Pick the lowest-index healthy channel; report none available.

\textbf{Hints.}
\begin{itemize}
\item Scan health flags; return the first healthy index, or $-1$.
\end{itemize}

\attempt

\textbf{Solution.}
\begin{cgym}
#include <stdio.h>

static int select_channel(const int *healthy, int n) {
    for (int i = 0; i < n; i++)
        if (healthy[i]) return i;
    return -1;                       /* all channels failed */
}

int main(void) {
    int h[3]  = {0, 1, 1};
    int hn[3] = {0, 0, 0};
    printf("%d %d\n", select_channel(h, 3), select_channel(hn, 3));   /* 1  -1 */
    return 0;
}
\end{cgym}

After voting flags a channel unhealthy, failover continues on the remaining
good channels; this picks a primary by priority order. The $-1$ ``no healthy
channel'' case is the one that matters most: the system must have a defined
safe response (revert to a backup law, hold last-good, or controlled
disengage) rather than using garbage --- the difference between fail-operational
and fail-safe (Q3.1). The health flags themselves come from the voter, parity/
CRC checks, and heartbeat/timeout monitors.

\followups
\begin{enumerate}
\item \emph{``Don't thrash between channels.''} $\rightarrow$ Add hysteresis ---
  require $N$ healthy frames before switching back.
\item \emph{``All channels fail --- then what?''} $\rightarrow$ A predefined
  safe state, never undefined behaviour; this is the critical path.
\item \emph{``Where do the health flags come from?''} $\rightarrow$ Voter
  disagreement, parity/CRC, SSM/status bits, and heartbeat timeouts.
\end{enumerate}
\end{gymproblem}

%=====================================================================
\section{Link to Her Resume}
%=====================================================================

\begin{resumelink}
\textbf{1553B / telemetry parsing:} ``At BEL I parsed MIL-STD-1553B and RS422
telemetry --- I decode command/status words and 429-style labels with explicit
shifts and masks, validate parity/CRC and the status/SSM before trusting any
value, and re-find the sync word on a bad frame rather than trusting byte
alignment.''

\medskip
\textbf{CCDL verification:} ``CCDL is the cross-channel link redundant flight
channels use to exchange data and vote each frame. I verify its integrity,
timing, and failover --- which means frame-synchronised channels, a
median/threshold vote that masks one bad channel without nuisance-tripping on
sensor noise, and a defined safe state when channels are lost.''

\medskip
\textbf{SPIL:} ``SPIL is a point-to-point serial processing link between cards;
I test it at the packet level --- framing, integrity, timing, and behaviour
under loss --- the same byte-level rigour I apply to every serial interface.''

\medskip
\small Maps directly to the resume: 1553B/RS422 telemetry parsing and the Python
automation at BEL, CCDL and SPIL verification at TASL. Keep claims to what the
resume states; the deep, defensible project scripting is Chapter 15, built only
from \code{inputs/INTAKE.md}.
\end{resumelink}

%=====================================================================
\section{Self-Test Scorecard}
%=====================================================================

\begin{scorecard}
\begin{enumerate}
\item What are the three 1553 terminal types and which initiates transfers?
\item Name the three 1553 word types and what each carries.
\item What fields are in an ARINC 429 word, and which one tells you if the data
  is valid?
\item Why is the ARINC 429 label commonly reversed in firmware?
\item Why does telemetry use UDP over TCP?
\item What does a CCDL carry, and why every control frame?
\item Why median (not average) voting, and why an odd number of channels?
\item Decode 1553 command word \code{0x2C22} (RT, T/R, SA, WC).
\item Why can't you cast a received avionics word onto a C bitfield struct?
\item Dual vs triple redundancy --- fail-safe or fail-operational, and why?
\end{enumerate}
\end{scorecard}

\begin{answerkey}
\begin{enumerate}
\item Bus Controller (initiates all transfers), Remote Terminal (responds when
  commanded), Bus Monitor (passive listener).
\item Command (BC$\rightarrow$RT: address/T-R/subaddr/word-count), status
  (RT$\rightarrow$BC: address + condition bits), data (16-bit payload).
\item Label, SDI, data, SSM, parity; the \textbf{SSM} reports validity (Normal /
  No-Computed-Data / Test / Failure).
\item The label is transmitted bit-reversed on the wire, so firmware reverses 8
  bits to recover the documented octal label.
\item A fresh sample supersedes a lost one; TCP's retransmission/ordering would
  stall the stream to re-deliver stale data.
\item Each channel's inputs/outputs to the others, every frame, so all channels
  vote on current data before driving the actuator.
\item Median rejects an outlier (an average is dragged toward it); an odd count
  gives a tie-breaker so one fault is outvoted.
\item RT 5, T/R = transmit (1), subaddress 1, word count 2.
\item Endianness, implementation-defined bitfield layout, and wire quirks
  (reversed label, packed fields) make struct overlay non-portable/corrupt ---
  use shifts and masks.
\item Dual = fail-safe (detect a disagreement but can't tell which is right);
  triple = fail-operational (outvote one fault and keep operating).
\end{enumerate}
\end{answerkey}

\begin{partnote}
\emph{End of Chapter 8, and of Part B (Communication Protocols).} Next: Part C
opens with \textbf{RTOS \& Real-Time} --- tasks and scheduling, priority
inversion, mutex vs semaphore, determinism, and the shared-memory/DPRAM handshake
patterns from the BEL work.
\end{partnote}
